\chapter{Движущееся ядро и его вторая фундаментальная форма}
\label{ch:v8-moving-kernel-second-fundamental-form-gate}

\section{Канонический оператор на страте полного ранга}

Для поля $X\in M_{3\times4}(\mathbb C)$ ранга три проектор на его ядерную
линию определяется без выбора базиса:
\begin{equation}
 P_X=I_4-X^*(XX^*)^{-1}X.
 \label{eq:v8-moving-kernel-projector}
\end{equation}
Его пространственно-временная производная является касательным вектором
грассманиана. Из $P_X^2=P_X$ следует точное разложение
\begin{equation}
 \nabla P_X=B_X+B_X^*,
 \qquad
 B_X=(I-P_X)(\nabla P_X)P_X.
 \label{eq:v8-moving-kernel-second-form}
\end{equation}
Блок $B_X$ переводит ядерную линию в трёхмерное дополнение и имеет ранг не
выше одного для каждого касательного направления.

Численный аудит подтвердил проекторные тождества и ковариантность при
$X\mapsto U_3XU_4^*$. Тем самым $B_X$ является настоящим каноническим
геометрическим объектом, а не новым произвольным отождествлением.

\section{Почему это ещё не недостающий коннектор}

На полном поле Тома VII естественный подъём имеет вид
\begin{equation}
 B_X\widehat\otimes I_{\mathcal Y_{\rm phys}}.
\end{equation}
Он коммутирует со всей алгеброй операторов физического рёберного множителя.
Следовательно, $B_X$ смешивает ядро и дополнение только внутри аффинного
$\mathbb C^4$, но не связывает рёберный и носители концевых секторов. Для постоянного
поля $\nabla P_X=0$, поэтому этот член даёт кинетическую геометрию семейной
ориентации, но не статический потенциал, массовую метрику или второй
семейный селектор.

\section{Стресс-тест настоящей смены ранга}

Рассмотрим семейство
\begin{equation}
 X_\varepsilon(x)=
 \begin{pmatrix}
  1&0&0&0\\
  0&1&0&0\\
  0&0&\varepsilon\cos(x/\varepsilon)&
       \varepsilon\sin(x/\varepsilon)
 \end{pmatrix}.
 \label{eq:v8-moving-kernel-singular-path}
\end{equation}
Норма $\partial_xX_\varepsilon(0)$ равна единице для всех
$\varepsilon>0$, однако
\begin{equation}
 \|B_{X_\varepsilon}(0)\|=\frac1{\varepsilon},
 \qquad
 \|B_{X_\varepsilon}(0)\|^2=\frac1{\varepsilon^2}.
\end{equation}
При $\varepsilon=0$ размерность ядра скачком меняется с одного до двух, а
операторная норма скачка проектора равна единице. Значит, стандартная
квадратичная норма второй фундаментальной формы не имеет универсального
непрерывного конечного продолжения через ранговую страту даже при
ограниченной производной исходного поля.

\begin{theorem}[Граница геометрии движущегося ядра]
На страте постоянного ранга $B_X$ каноничен, базисно ковариантен и задаёт
естественную грассманианную кинетическую метрику ядерной линии. Но он
факторизован относительно физического рёберного носителя, исчезает на
постоянных конфигурациях и может расходиться как $1/\varepsilon$ при
приближении к смене ранга. Поэтому он не является недостающим
межсекторным родителем Тома VIII.
\end{theorem}

\begin{nogocriterion}[Запрет объявлять $B_X$ регуляризатором ранга]
Нельзя считать проекторную производную автоматически гладкой в точке
изменения ранга. Новый вариант обязан либо вывести нефакторизованный
оператор, либо задать независимый регуляризованный проектор и доказать
конечность действия без ручного обрезание по минимальному сингулярному числу.
\end{nogocriterion}

\begin{protocol}
\begin{enumerate}
 \item канонический проектор ядра на ранга три страте: \textbf{получен};
 \item вторая фундаментальная форма: \textbf{получена};
 \item базисная и калибровочная ковариантность: \textbf{пройдена};
 \item связь ядра с аффинным дополнением: \textbf{есть, ранг не выше одного};
 \item связь рёберного и концевом носителей: \textbf{нет};
 \item статическая массовая метрика: \textbf{не возникает};
 \item конечность на общем изменяющимся рангом пути: \textbf{провалена};
 \item внутренний поиск основания Тома VIII: \textbf{заморожен}.
\end{enumerate}
\end{protocol}

Машинный сертификат сохранён в
\path{s2t_v8_moving_kernel_second_fundamental_form_gate_results.json}.